\documentclass[10pt,a4paper]{article}

\usepackage{iftex}
\ifPDFTeX
  \usepackage[utf8]{inputenc}
  \usepackage[T1]{fontenc}
  \usepackage{lmodern}
\else
  \usepackage{fontspec}
\fi

\usepackage[magyar]{babel}
\usepackage[a4paper,margin=1.75cm,top=1.9cm,bottom=1.85cm]{geometry}
\usepackage{xcolor}
\usepackage{booktabs}
\usepackage{tabularx}
\usepackage{enumitem}
\usepackage{fancyhdr}
\usepackage{titlesec}
\usepackage{hyperref}

% Spera arculati szinek (eszaki teli paletta nyomtatasbarat valtozata)
\definecolor{speraSand}{HTML}{EDF3F7}
\definecolor{speraTerracotta}{HTML}{1F637A}
\definecolor{speraPeach}{HTML}{B88428}
\definecolor{speraSage}{HTML}{2B6E55}
\definecolor{speraInk}{HTML}{14202E}
\definecolor{speraMuted}{HTML}{4C5E70}

\hypersetup{
  colorlinks=true,
  linkcolor=speraTerracotta,
  urlcolor=speraTerracotta,
  pdftitle={Spera -- FIESZTA 2026 Dokumentacio},
  pdfauthor={Jelige: Manci}
}

\setlength{\parindent}{0pt}
\setlength{\parskip}{4pt}
\setlist[itemize]{leftmargin=1.2em,itemsep=2pt,topsep=2pt,parsep=0pt}
\setlist[enumerate]{leftmargin=1.4em,itemsep=2pt,topsep=2pt,parsep=0pt}

\titleformat{\section}
  {\large\bfseries\color{speraTerracotta}}
  {\thesection.}{0.5em}{}[\vspace{-4pt}\textcolor{speraPeach}{\rule{\linewidth}{0.8pt}}]
\titlespacing*{\section}{0pt}{9pt}{4pt}

\titleformat{\subsection}
  {\normalsize\bfseries\color{speraInk}}
  {\thesubsection}{0.5em}{}
\titlespacing*{\subsection}{0pt}{6pt}{2pt}

\pagestyle{fancy}
\fancyhf{}
\fancyhead[L]{\small\color{speraMuted}Spera -- Nőgyógyászati-endokrinológiai személyes egészségügyi dashboard}
\fancyhead[R]{\small\color{speraMuted}FIESZTA 2026 | Jelige: \textbf{Manci}}
\fancyfoot[C]{\small\color{speraMuted}\thepage}
\renewcommand{\headrulewidth}{0.4pt}
\renewcommand{\footrulewidth}{0pt}

\begin{document}

% Cimfejlec es metaadat-sav
\noindent\colorbox{speraSand}{%
  \parbox{\dimexpr\linewidth-2\fboxsep\relax}{%
    \vspace{4pt}
    {\Large\bfseries\color{speraTerracotta} Spera} \hfill
    {\small\bfseries\color{speraSage} FIESZTA 2026 Pályázat \,|\, Jelige: Manci}\\[3pt]
    {\normalsize\bfseries\color{speraInk} Menstruációs cikluskövetéssel integrált, nőgyógyászati-endokrinológiai személyes egészségügyi dashboard}\\[4pt]
    {\small\color{speraMuted} Kísérő dokumentáció, betegút-elemzés, indítási útmutató és üzleti-fejlesztési összefoglaló (max. 5 oldal).}
    \vspace{4pt}
  }%
}

\vspace{2pt}

\section{Bevezetés, választott kórképek és fiktív célfelhasználó}

A \textbf{Spera} egy 100\%-ban offline működő, betegközpontú személyes egészségügyi dashboard, amely a menstruációs cikluskövetést kapcsolja össze a krónikus nőgyógyászati-endokrinológiai állapotok mindennapi önmenedzsmentjével. A rendszer modulárisan, egymással kombinálhatóan támogatja a reproduktív korú nőket leginkább érintő három kórképcsoportot -- valamint az általános cikluskövetést és kivizsgálás alatti állapotot:
\begin{itemize}
  \item \textbf{PCOS / Inzulinrezisztencia (IR)} (\texttt{E28.2 / E11.9}): ciklusszabályosság, eltolódott ovuláció, HOMA-IR és vércukor-ingadozás követése.
  \item \textbf{Hashimoto-thyreoiditis / pajzsmirigy-alulműködés} (\texttt{E06.3 / E03.9}): TSH és fT4 trendek, L-tiroxin időzítés és mikrotápanyag-kölcsönhatások figyelése.
  \item \textbf{Endometriózis / ösztrogéndominancia} (\texttt{N80}): kismedencei fájdalom, vérzésintenzitás, gyulladáscsökkentő étrend és vasraktárak követése.
\end{itemize}

\textbf{Fiktív célfelhasználó (persona):} \textit{Szabó Eszter} (29 éves, vércsoport: A Rh+, allergiák: penicillin -- súlyos, laktóz -- enyhe). 2022-ben PCOS/IR, 2023-ban emelkedett anti-TPO (215~IU/ml) alapján Hashimoto-thyreoiditis diagnózist kapott. Jelenlegi terápiája: \textit{Letrox 50~$\mu$g} (06:30, éhgyomorra), \textit{Merckformin XR 500~mg} (08:00 és 20:00), \textit{Mio-inozitol + D3-vitamin 2000~NE} (13:00), \textit{Magnézium-biszglicinát 300~mg} (21:30). Bár az alapértelmezett ciklushossza 28~nap, az elmúlt három ciklusa 33--35~naposra nyúlt, miközben 5~nap múlva esedékes endokrinológiai és nőgyógyászati kontrollvizite előtt áll, és babatervezést fontolgat.

\section{A megoldandó probléma és a dashboard célja}

Az endokrin és nőgyógyászati krónikus betegséggel élő nők egészségügyi adatai a gyakorlatban széttagoltak: a vérzésnapló egy általános ciklusalkalmazásban van, a hormonlaborok PDF-ként az EESZT-ben, az ultrahang-leletek papíron, míg a gyógyszerek és mikrotápanyagok bevételi szabályai fejben maradnak. Ebből négy kritikus probléma fakad:
\begin{enumerate}
  \item \textbf{Kontextus nélküli tünetek:} az általános ciklusappok nem kezelik az endokrin kórképek fázisfüggő tüneteit (pl. luteális fázisban felerősödő pajzsmirigy-fáradtság, étkezés utáni álmosság).
  \item \textbf{Rosszul időzített hormonlaborok:} egy 34~napos ciklusban a sematikusan a 21. napra kért progeszteronvétel még az ovuláció előtt történik, így tévesen alacsony eredményt ad.
  \item \textbf{Gyógyszer--mikrotápanyag ütközések:} az L-tiroxin felszívódását a túl korán bevitt vas, kalcium vagy kávé jelentősen rontja.
  \item \textbf{Szűkös vizitidő:} a 10--15 perces szakorvosi konzultáción a beteg gyakran nem tudja pontosan felidézni az elmúlt hónapok ciklushosszát, gyógyszerhűségét (adherence) és kérdéseit.
\end{enumerate}

A \textbf{Spera célja}, hogy mindezeket egyetlen, nyugodt és átlátható felületen egyesítse, elősegítve az adatok értelmezését, a helyes labor- és gyógyszer-időzítést, valamint a felkészült orvos--beteg kommunikációt.

\section{Fő funkciók, technológiai felépítés és indítási útmutató}

\subsection{A prototípus főbb funkciói (4 redundanciamentes, mobilbarát nézet)}
\begin{itemize}
  \item \textbf{Áttekintő:} 6 hetes (1,5 hónapos) gördülő, mind a 4 ciklusfázist színkóddal és ciklusnappal jelölő naptár, napra pontos menstruáció-előrejelzés, \textit{Smart Alert} (automatikus ciklushossz-frissítés), ciklusnap-függő hormonlabor-emlékeztető, \textbf{Mai gyógyszerek} gyors-pipálás és az aktuális fázishoz igazított \textbf{Étrend \& Vitamin ajánló} kölcsönhatás-figyelővel.
  \item \textbf{Egykattintásos Élethelyzet-váltó és akadálymentesség:} a fejlécben bármikor váltható a \textit{Cikluskövetés}, \textit{Baba tervezése} (prekoncepciós orvosi ellenőrzőlista, LH-teszt) és \textit{Várandósság követése} (heti gyümölcs/zöldség méret-összehasonlítás, vérnyomás-, vércukor- és magzatmozgás-naplózás, kötelező szűrési mérföldkövek) mód, valamint elérhető az \texttt{A+} betűméret-kapcsoló. A felület vízszintes görgetés nélkül, mobilon fix alsó navigációs sávval használható.
  \item \textbf{Gyorsteszt \& Napló:} egységesített, nőgyógyászati-endokrinológiai fókuszú napi állapotfelmérés (gyógyszerhűség, energia, agyköd, alvás, hangulat, ciklus-biomarkerek, BBT, valamint beépített metabolikus/pajzsmirigy/androgén klinikai jelek) és vizit előtti felkészítő űrlap kikapcsolható \textit{Orvosi háttér} magyarázatokkal.
  \item \textbf{Elemzések \& Laborok:} 30 napos önmenedzsment-mintázatok összesítő és interaktív, inline SVG grafikonok -- ciklushossz-trend (21--35 napos referencia-sávval), szűrhető endokrin labor-grafikon normál és optimális célérték-sávokkal, valamint BBT- és tünetintenzitás-görbe.
  \item \textbf{Betegút \& Orvosi adatok:} interaktív 5 állomásos betegúttérkép, közvetlenül alatta funkcióként az 1 kattintásos, nyomtatható A4 \textbf{Orvosi összefoglaló (Állapotlap)}, valamint minden fontos orvosi adat (kórképek BNO-kóddal és fázissal, gyógyszerek, allergiák, kórelőzmény, JSON mentés).
\end{itemize}

\subsection{Technológiai felépítés, indítási útmutató és ismert korlátozások}
A prototípus külső keretrendszer, webszerver, build lépés, adatbázis és internetkapcsolat nélkül, tiszta HTML5, CSS3 és ES6+ Vanilla JavaScript (\texttt{index.html}, \texttt{styles.css}, \texttt{data.js}, \texttt{app.js}) technológiával készült. Minden grafikon inline SVG-ként renderelődik, a stílusok és a favicon beágyazottak, nincs \texttt{fetch()} hívás vagy külső CDN, így \texttt{file://} protokollon CORS-hiba nélkül fut. A felhasználói adatok kizárólag helyben, a böngésző \texttt{localStorage} tárhelyén (\texttt{spera\_fieszta\_2026\_v1}) tárolódnak, minden dinamikus DOM-kiírás XSS-védett (\texttt{escapeHtml()}).

\textbf{Indítási útmutató:}
\begin{enumerate}
  \item Csomagold ki a \texttt{spera\_Manci\_prototipus.zip} állományt egy tetszőleges helyi mappába.
  \item Nyisd meg az \texttt{index.html} fájlt dupla kattintással Chrome, Edge vagy Firefox böngészőben (\texttt{file://} protokoll, internet nélkül is működik).
  \item Az induló kérdőív kitölthető saját adatokkal, vagy a képernyő alján található \textit{„Demó profil automatikus kitöltése”} gombbal azonnal betölthető \textit{Szabó Eszter} teljes szintetikus betegprofilja.
\end{enumerate}

\textbf{Ismert technikai korlátozások:} a helyi \texttt{localStorage} böngészőhöz és eszközhöz kötött (eszközök közötti felhőszinkron nélkül); a demóadatok teljes mértékben szintetikusak; a prototípus tájékoztató és önmenedzsment-támogató eszköz, nem minősül orvostechnikai eszköznek és nem helyettesíti a szakorvosi ellátást.

\section{Szintetikus adatok és orvosszakmai források}

A \texttt{data.js} modul (\texttt{window.SPERA\_DATA}) tartalmazza a szintetikus betegprofilt és a szakmai tudásbázist. A \texttt{shiftDate()} segédfüggvény a mindenkori aktuális dátumhoz viszonyítva generálja a ciklusnapokat, a vizit előtti 5 napos visszaszámlálást és a laboridőpontokat, így a demó bármikor megnyitva naprakész. A 7 követett biomarker (\texttt{TSH}, \texttt{fT4}, \texttt{HOMA-IR}, \texttt{Progeszteron}, \texttt{Ösztradiol}, \texttt{D-vitamin}, \texttt{Ferritin}) normál és optimális/babatervezési célsávjai az alábbi nemzetközi irányelvekre épülnek:
\begin{itemize}
  \item \textbf{PCOS / IR:} \textit{International Evidence-based Guideline for the Assessment and Management of PCOS (ESHRE/ASRM, 2023)} -- cikluszavarok, HOMA-IR, mio-inozitol és életmód-terápia.
  \item \textbf{Pajzsmirigy és várandósság:} \textit{American Thyroid Association (ATA, 2017) Guidelines for Thyroid Disease During Pregnancy and the Postpartum} -- babatervezési TSH céltartomány ($<2{,}5$~mIU/l), L-tiroxin felszívódási interakciók (vas, kalcium, kávé).
  \item \textbf{Endometriózis:} \textit{ESHRE Endometriosis Guideline (2022)} -- fájdalomnaplózás, gyulladáscsökkentő étrendi támogatás és vaspótlás erős vérzés esetén.
\end{itemize}

\section{A betegút-ábrázolás magyarázata}

A prototípus \textit{Betegúttérkép} nézete öt, egymásra épülő ellátási állomáson keresztül mutatja be a nőgyógyászati-endokrinológiai gondozást. Minden állomáshoz tartozik a beteg aktuális döntési helyzete, a magyar ellátórendszerben gyakori elakadás, valamint a dashboard konkrét támogatási pontja és közvetlen navigációs gombja:

\begin{center}
\small
\renewcommand{\arraystretch}{1.18}
\begin{tabularx}{\linewidth}{@{}p{3.1cm} X X X@{}}
\toprule
\textbf{Ellátási állomás} & \textbf{Beteg döntési helyzete} & \textbf{Tipikus elakadás} & \textbf{Spera támogatási pont} \\
\midrule
\textbf{1. Tünetfelismerés és önmegfigyelés} & Normális ciklus-ingadozás vagy endokrin kivizsgálást igénylő jel? & Tünetek szétaprózottsága, a ciklusfázis és a panaszok közti kapcsolat fel nem ismerése. & Fázis- és kórképspecifikus tünetnapló, BBT- és metabolikus mikrotünet-követés. \\
\textbf{2. Első ellátási pont} (háziorvos / nőgyógyász) & Elegendő-e a tüneti kezelés, vagy szakorvosi beutalás indokolt? & Panaszok elbagatellizálása, kivizsgálás nélküli terápia, viziten elfelejtett adatok. & 1 kattintásos, nyomtatható Orvosi összefoglaló konzultációs kérdésekkel és adherence-aránnyal. \\
\textbf{3. Célzott endokrin diagnosztika} & Melyik ciklusnapon történjen a vérvétel az értékelhető eredményhez? & Hosszabb ciklusban a 21. napon levett progeszteron tévesen alacsony értéket mutat. & Smart Alert a valódi ciklushosszhoz, 3--5. napi és ovuláció + 7. napi labor-emlékeztető. \\
\textbf{4. Diagnózis és személyre szabott terápia} & Hogyan illeszthetők össze biztonságosan a gyógyszerek és vitaminok? & Bonyolult bevételi szabályok (L-tiroxin vs. vas/kávé), kimaradó adagok miatti visszaesés. & Kórkép-adattár, időzített gyógyszer-emlékeztető és automatikus kölcsönhatás-figyelő. \\
\textbf{5. Hosszú távú gondozás és családtervezés} & Hogyan változnak a célértékek (pl. TSH) babatervezés vagy terhesség alatt? & Pozitív teszt után elkésett TSH-kontroll vagy elmaradó korai OGTT szűrés. & Egykattintásos Élethelyzet-váltó és azonnali trimeszter-specifikus endokrin riasztások. \\
\bottomrule
\end{tabularx}
\end{center}

\section{Üzleti, megvalósíthatósági és továbbfejlesztési koncepció}

\begin{itemize}
  \item \textbf{Célcsoport és értékajánlat:} Magyarországon a reproduktív korú nők mintegy 10--15\%-a érintett PCOS-ben, pajzsmirigy-alulműködésben vagy endometriózisban. A Spera a betegek számára csökkenti a mentális terhet és javítja a terápiahűséget, az endokrinológusok és nőgyógyászok számára pedig strukturált, azonnal áttekinthető anamnézist ad a vizitek elején.
  \item \textbf{Továbbfejlesztés és rendszerkapcsolatok:} A jelenlegi offline váz moduláris adatmodellje (\texttt{window.SPERA\_DATA} és normalizált állapot) közvetlenül alkalmas \textbf{EESZT / HL7 FHIR} (\texttt{Observation}, \texttt{MedicationStatement}, \texttt{Condition}) alapú automatikus laborimportra, hordozható eszközökről (okosóra, BBT-gyűrű) érkező hőmérséklet- és alvásadatok szinkronizálására, valamint PDF/QR-kódos leletátadásra.
  \item \textbf{Adatvédelmi és szabályozási feltételek:} Egészségügyi adatok esetén (GDPR 9. cikk szerinti különleges adatok) a \textit{local-first} architektúra minimális adatkezelési kockázatot jelent; későbbi felhőszinkron esetén végponti titkosítás és explicit beteg-hozzájárulás szükséges. Klinikai döntéstámogató funkciók bővítésekor az EU MDR (2017/745) szoftver-osztályozási (SaMD) követelményei irányadók.
  \item \textbf{Fenntarthatósági és finanszírozási modell:} Kétlépcsős, hibrid modell: (1) \textbf{B2C Freemium} -- az alapvető cikluskövetés, tünetnapló és betegút ingyenes, míg a fejlett labor-trend elemzés és export előfizetéses; (2) \textbf{B2B2C magánklinikai licenc} -- nőgyógyászati-endokrinológiai centrumok számára fehér címkés vizit-előkészítő modulként, csökkentve az adminisztrációs vizitidőt.
\end{itemize}

\end{document}
